\chapter*{Introduzione}
\addcontentsline{toc}{chapter}{Introduzione}

\section*{Contesto e motivazioni}

Negli ultimi decenni il Machine Learning (ML) si è affermato come lo strumento computazionale di riferimento per l'estrazione di pattern predittivi da dati complessi, trovando applicazione in ambiti eterogenei quali la diagnostica medica, l'analisi chimica e il riconoscimento di immagini\footnote{Mitchell, T. M. (1997). \textit{Machine Learning}. McGraw-Hill, New York.}. Parallelamente, il calcolo quantistico ha compiuto negli ultimi anni progressi significativi, passando da un ambito prevalentemente teorico alla disponibilità di dispositivi fisici accessibili tramite piattaforme cloud, seppure ancora vincolati dalle limitazioni proprie dell'attuale era NISQ (\textit{Noisy Intermediate-Scale Quantum})\footnote{Preskill, J. (2018). Quantum Computing in the NISQ era and beyond. \textit{Quantum}, 2, 79.}.
All'intersezione di queste due discipline si colloca il Quantum Machine Learning (QML), un campo di ricerca emergente che propone di sfruttare le proprietà della meccanica quantistica, quali la sovrapposizione e l'entanglement, per rappresentare ed elaborare l'informazione in spazi di Hilbert la cui dimensione cresce esponenzialmente con il numero di qubit impiegati\footnote{Biamonte, J., Wittek, P., Pancotti, N., Rebentrost, P., Wiebe, N., \& Lloyd, S. (2017). Quantum machine learning. \textit{Nature}, 549(7671), 195-202.}. Questa capacità rappresentazionale ha motivato l'introduzione di modelli di classificazione quantistica, quali il Quantum Support Vector Classifier (QSVC) e il Variational Quantum Classifier (VQC), proposti in letteratura come potenziali alternative ai classificatori classici, in particolare per problemi caratterizzati da strutture di correlazione tra le feature difficili da catturare efficientemente con un kernel classico\footnote{Havlíček, V., Córcoles, A. D., Temme, K., Harrow, A. W., Kandala, A., Chow, J. M., \& Gambetta, J. M. (2019). Supervised learning with quantum-enhanced feature spaces. \textit{Nature}, 567(7747), 209--212.}.
Tuttavia gli articoli scientifici che trattano questo argomento sono ancora per lo più studi puramente teorici oppure esperimenti empirici condotti su set di dati fabbricati esplicitamente per dimostrare l'avvantaggio dei modelli quantistici rispetto a quelli classici mentre vi è solo poca evidenza empirica sull'utilizzo di tali algoritmi su set di dati reali che sono comuni nella pratica del Machine Learning classico. L'obiettivo della tesi è quello di cercare di colmare almeno parzialmente questo divario.

\section*{Obiettivi della tesi}

L'obiettivo generale di questo lavoro è valutare, in modo empirico e metodologicamente rigoroso, se ed in quale misura i modelli di Quantum Machine Learning siano in grado di competere con i modelli di classificazione classici consolidati, quando applicati a dataset reali sotto i vincoli computazionali attualmente imposti dall'hardware quantistico disponibile. Da questo obiettivo generale derivano i seguenti obiettivi specifici:

\begin{itemize}
    \item formalizzare il quadro teorico e matematico necessario al confronto, sia per la classificazione classica (paradigmi di apprendimento, algoritmi di riferimento, metodologia di valutazione) sia per il Quantum Machine Learning (data encoding, circuiti quantistici parametrizzati, architetture QSVC e VQC);
    \item progettare e implementare un sistema sperimentale modulare e riproducibile, in grado di addestrare e valutare, a parità di dati in ingresso e di protocollo di valutazione, sia i modelli classici sia i modelli quantistici;
    \item condurre una campagna sperimentale su tre dataset reali, scelti per rappresentare condizioni di classificazione binaria e multiclasse a diversa dimensionalità e provenienti da domini applicativi distinti;
    \item confrontare le prestazioni predittive e il costo computazionale dei due paradigmi, sia in condizioni di simulazione ideale sia in presenza di un modello di rumore realistico, discutendo criticamente i risultati alla luce dei limiti teorici e pratici del QML nell'attuale era NISQ.
\end{itemize}

\section*{Metodologia adottata}

Per raggiungere questi obiettivi l'approccio utilizza un protocollo sperimentale comune utilizzato per tutti gli algoritmi confrontati: tre algoritmi benchmark classici (Regressione Logistica, SVM con kernel RBF,Random Forest), per i quali i parametri iperparametri sono ottimizzati utilizzando la ricerca a griglia con la validazione incrociata e due architetture quantistiche create utilizzando il framework Qiskit (QSVC,VQC), entrambi applicati a tre diversi set di dati reali disponibili tramite Scikit-learn (Breast Cancer Wisconsin, Wine e Digits).
I set di dati sono normalizzati e, se necessario, la loro dimensione viene ridotta tramite PCA per renderli compatibili con il piccolo numero di qubit che possono essere gestiti con la simulazione. L'intero ambiente è containerizzato utilizzando Docker per garantire la precisa replicabilità degli esperimenti. La metodologia e l'implementazione di questo protocollo sono descritte rispettivamente nei capitoli~\ref{cap:architettura_sistema_metodologia} e~\ref{cap:implementazione_sperimentale_risultati}.

\section*{Struttura della tesi}

I capitoli rimanenti di questa tesi sono strutturati nel seguente modo. Il Capitolo~\ref{cap:fondamenti_ml_classificazione_classica} (Fondamenti del Machine Learning e della Classificazione Classica) introduce i paradigmi del Machine Learning, la formulazione matematica del problema della classificazione supervisionata, i metodi di pre-elaborazione dei dati reali e gli algoritmi classici utilizzati come linea di base, nonché l'approccio alla valutazione delle loro prestazioni.
Il Capitolo~\ref{cap:quantum_machine_learning} (Machine Learning Quantistico) fornisce la definizione formale dei fondamenti dell'informatica quantistica nel contesto del Machine Learning, i metodi di codifica dei dati e l'architettura dei classificatori QSVC e VQC, nonché discutendo i limiti dell'era NISQ, come i piatti deserti e il rumore hardware.
Il Capitolo~\ref{cap:architettura_sistema_metodologia} (Architettura del Sistema e Metodologia) spiega l'architettura del sistema sperimentale creato, lo stack di tecnologie utilizzate, il dataset scelto e la metodologia sperimentale seguita per le comparazioni.
Il Capitolo~\ref{cap:implementazione_sperimentale_risultati} (Implementazione Sperimentale e Risultati) descrive l'implementazione software del sistema sperimentale e mostra i risultati ottenuti dagli esperimenti condotti sui tre dataset, confrontando modelli classici e quantistici e analizzando criticamente i loro svantaggi.
